% ============================================================================
% Appendix A9 — eps-Extrapolation of Equilibrium Quantities
% ============================================================================
\section{$\varepsilon$-Extrapolation of Equilibrium Quantities}
\label{app:epsilon}

Analytical statements in this paper are made for the regularized system
($\varepsilon>0$) with the $\varepsilon\to0$ limit available in closed form
(Appendix~\ref{app:equilibrium}); the companion code integrates at
$\varepsilon=0.1$. Because the scope-absorption invariant
$\beta C^{*\eta}=O^{*}-O_0$ (Corollary~\ref{cor:scope}) is itself
$\varepsilon$-dependent, quantitative uses of the invariant must quote the
regularization at which it was computed. Table~\ref{tab:epsilon} reports
the interior equilibrium across a decreasing $\varepsilon$ ladder, computed
with the production root-finder
(\texttt{find\_interior\_equilibria}, $C$-window $200$, $16{,}000$ scan
points), together with the exact $\varepsilon\to0$ closed form of
Eq.~\eqref{eq:eps0-quadratic}.

\begin{table}[h]
\centering
\caption{Interior equilibrium versus the barrier regularization
$\varepsilon$ at baseline parameters. The bottom row is the closed-form
$\varepsilon\to0$ limit, not an extrapolation.}
\label{tab:epsilon}
\vspace{0.4em}
\begin{tabular}{lcccc}
\toprule
$\varepsilon$ & $V^{*}$ & $C^{*}$ & $O^{*}$ & $\beta C^{*}$ \\
\midrule
$0.2$   & $4.72457$ & $32.32766$ & $9.08191$ & $8.08191$ \\
$0.1$ (baseline) & $4.70253$ & $32.03372$ & $9.00843$ & $8.00843$ \\
$0.05$  & $4.69130$ & $31.88406$ & $8.97102$ & $7.97102$ \\
$0.02$  & $4.68450$ & $31.79338$ & $8.94834$ & $7.94834$ \\
$0.01$  & $4.68222$ & $31.76300$ & $8.94075$ & $7.94075$ \\
\midrule
$0$ (closed form) & $4.67994$ & $31.73254$ & $8.93313$ & $7.93313$ \\
\bottomrule
\end{tabular}
\end{table}

Three observations.

\begin{enumerate}
\item \emph{Monotone, small, first-order.} Every quantity approaches its
limit monotonically from above, with increments approximately proportional
to $\varepsilon$ (halving $\varepsilon$ roughly halves the remaining gap).
The total displacement from the working value $\varepsilon=0.1$ to the
limit is $\Delta V^{*}=0.0226$ ($0.48\%$) and
$\Delta(\beta C^{*})=0.0753$ ($0.94\%$)---the source of the ``$<2\%$
shift'' statement in the companion code documentation, here sharpened.
\item \emph{The invariant's two quotable values.} The scope-absorption
product is $\beta C^{*}=8.008$ at the working regularization
$\varepsilon=0.1$ and exactly $7.933$ ($=O^{*}-O_0$ from the quadratic
\eqref{eq:eps0-quadratic}) in the $\varepsilon\to0$ closed form. Both are
correct; neither should be quoted without its $\varepsilon$. The pinning
tests use the $\varepsilon=0.1$ value.
\item \emph{Classifications are $\varepsilon$-robust here.} The regime
label at baseline (low-vitality: $V^{*}<\Vstrat=5$) is unchanged across
the entire ladder, and the regime flip value of $\mu/\alpha$ moves only
from $2.163$ ($\varepsilon=0.1$) to $2.175$ ($\varepsilon\to0$)
(Appendix~\ref{app:equilibrium}). For calibrations nearer a regime
boundary than the baseline's $8\%$ margin, the $\varepsilon$-ladder should
be re-run before any classification is reported---this is the practical
protocol the table exists to demonstrate.
\end{enumerate}
